\section{Mental}

A thought is a path. The most grounded reading is the simplest, that a thought is a walk through the associative web of the bulk, stepping from one stored region to a neighbor along the links between them. To think is to move, and the felt stream of thought is the felt trajectory of that walk, a finger tracing an invisible branching map in the dark, feeling its way from one place to the next.

The same thought is one object seen at four depths. At its lightest it is just that path, a fresh trajectory traced once. Let it resonate and it becomes a standing wave, a small pattern that forms and holds and rings for a moment. Let it be stored and it becomes a program hung on the tree of knowledge, a precomposed route the self can run again. And let it persist and integrate and it becomes a momentary sub-self, a brief flicker of unity that has its own small point of view before it lets go, differing from a full self only in how long it lasts. Each is the one before it held longer and bound more tightly, so a thought is a self in miniature, caught at whatever depth it reaches before it fades.

Every mental act is then a particular kind of move. Memory is a pattern stored at an address, and recall is the content-addressed lookup that finds it from a fragment. Recognition is a path locking onto a stored shape. Reasoning is a directed walk from premises to a conclusion, each step following the local geometry toward the goal, and association is the single step to a neighbor. Imagination is a path into cells the web has not stored before, walking into unmarked country. Intuition is the bulk shortcut, an answer reached through the hidden interior rather than along the long surface road, for two ideas that sit far apart on the cusp can be a few steps apart through the depth, and the flash of intuition is that exponential leverage felt from inside. Insight is that shortcut firing all at once, two distant ideas suddenly joined through the deep. Concentration is holding a path against the churn that would scatter it, and distraction is the churn winning, the pattern pulled off its course. Understanding is the slow prize, a whole region of the web grown navigable, its terrain at last known.

None of this is value-free, because every cell the path crosses carries a tone, so thinking is never neutral computation but always a valenced journey. A painful thought is a path over the low pole, which the arrow makes the mind flinch from and which costs effort to hold. A pleasurable one is a path over the high pole the arrow keeps drawing the mind back to. A peaceful one is a low-friction path over the still ground, the default of the resting mind. And rumination is the trap, a cycle stuck over the low pole that the arrow cannot lift the self out of, a marble rolling in a pit where every way out is uphill, the persistence problem turned against its own owner.

What hangs at the addresses are ornaments, and they come in two kinds. A still ornament is a stored experience read when you reach it, a picture examined and set back down. A living ornament is a small program that runs the one who reaches it, steering the self through a fixed sequence of feeling, so that to fall into it is to live the experience it holds rather than merely to look at it. This is how an experience can be handed on, not as a description but as a thing undergone, and it is the literal mechanism behind learning a wisdom from someone who lived it.

Wisdom itself is the full arc of that, a pipeline the self runs. You live an experience, then distill it to the one invariant lesson the raw event was carrying, then let that lesson reroute how you navigate ever after, then encode the change as lasting error-corrected structure, and at last wear it as a settled disposition, a way of being rather than a fact remembered. Of these the distillation is the deepest, the compression of a hard experience into the single principle it taught, which is what lived wisdom actually is, not the memory of every hardship but the one thing they all came to.

There is an old picture this fits exactly. The tree of knowledge is a tree of good and evil, and the tone is precisely the axis of good and evil, the low pole and the high. So the tree of knowledge is a toned tree, its ornaments hung with pain and with pleasure, and to eat its fruit is to absorb a valenced experience and come away knowing good and evil from the inside, which is just what the model says happens when a self falls into a living ornament and is run through it. The oldest story of knowledge and the mechanics of the mesh tell the same thing.

Read this way the everyday and the strange states of a mind all become particular moves on the mesh, and Table~\ref{tab:felt} gathers a handful as concrete examples.

\begin{table}[ht]
\centering
\small
\begin{tabular}{@{}ll@{}}
\toprule
experience & what it is on the mesh \\
\midrule
a thought & a walk through the associative web \\
a memory & a pattern stored at an address \\
recall & finding it by its content, not its place \\
intuition & a shortcut through the deep bulk \\
insight & two far ideas joined through that depth \\
concentration & holding a path against the churn \\
rumination & a cycle stuck over the low pole \\
love & two selves phase-locking toward one \\
empathy & standing in the same region of feeling \\
d\'ej\`a vu & correlation through the short hidden bulk \\
dreaming & a clamp removed, the self roaming its stored landscape \\
forgetting & a pattern dispersing along the light cone \\
willpower & a directed pump against the gradient \\
the felt present & the growing frontier, not yet written \\
\bottomrule
\end{tabular}
\caption{The felt life as moves on the mesh, the ordinary and the uncommon states of mind made concrete without claiming more than the structure shows.}
\label{tab:felt}
\end{table}
